\section{A fixed finite prime family detects every boundary monomial}

\begin{definition}[The fixed boundary primes]\label{def:boundary-minimal-primes}
\lean{EnvelopeProperty.boundaryMinimalPrimes}
\leanok
For a finite family $p_i\in E$, take the union, over all $i$, of the
primes minimal over $(p_i)$. This family is independent of every later
choice of monomial exponents.
\end{definition}

\begin{lemma}[Finiteness of the actual fixed prime family]\label{lem:finite-boundary-minimal-primes}
\lean{EnvelopeProperty.finite_boundaryMinimalPrimes}
\leanok
\uses{def:boundary-minimal-primes}
If $E$ is Noetherian, the fixed boundary prime family is finite.
\end{lemma}
\begin{proof}
Every ideal in a Noetherian ring has only finitely many minimal primes.
The defining union has a finite index set, so is a finite union of
finite sets.
\end{proof}

\begin{lemma}[Every monomial minimal prime is a fixed boundary prime]\label{lem:monomial-minimal-primes-boundary}
\lean{EnvelopeProperty.minimalPrimes_prod_pow_subset_boundaryMinimalPrimes}
\leanok
\uses{def:boundary-minimal-primes}
For arbitrary natural exponents $e_i$, every prime minimal over
$(\prod_i p_i^{e_i})$ is minimal over $(p_i)$ for some $i$ with $e_i>0$.
\end{lemma}
\begin{proof}
Write $g=\prod_i p_i^{e_i}$ and take a prime $P$ minimal over $(g)$.
Since $g\in P$ and $P$ is prime, some active factor has $p_i\in P$
and $e_i>0$. Otherwise every factor lies in the prime complement and
so does the product, a contradiction. If $(p_i)\subseteq Q\subseteq P$
with $Q$ prime, then $p_i\mid g$ makes $g\in Q$. Minimality of $P$ over
$(g)$ gives $P\subseteq Q$, hence $P=Q$. Thus $P$ is minimal over $(p_i)$.
\end{proof}

\begin{theorem}[Fixed-family localization detection for every monomial]\label{thm:fixed-boundary-localization-detection}
\lean{EnvelopeProperty.mem_span_prod_pow_iff_boundary_localizations}
\leanok
\uses{lem:monomial-minimal-primes-boundary,thm:regular-principal-dvr-detection,lem:finite-boundary-minimal-primes}
Let $E$ be a regular domain and let every $p_i$ be nonzero. For arbitrary
exponents $e_i$, an element $f\in E$ belongs to $(\prod_i p_i^{e_i})$
if and only if it belongs to that principal ideal after localization
at every prime in the fixed boundary family. Every such localization
is a DVR, and the finite testing family is independent of $e$.
\end{theorem}
\begin{proof}
Global divisibility persists in every localization. Conversely the
product is nonzero. The regular-domain principal-ideal criterion tests
membership using the actual DVR localizations at primes minimal over
its principal ideal. The preceding lemma puts every one of those
primes in the fixed boundary family, so the assumed tests imply
membership in the original ring. Each boundary prime is minimal over
a nonzero $(p_i)$; the same regular-domain theorem makes its localization
a DVR. Finiteness follows from Noetherianity of the regular domain.
\end{proof}
